\documentclass{article}
\usepackage[utf8]{inputenc}
\usepackage{amsmath}
\usepackage{amssymb}

\title{核心数学模型}
\date{}

\usepackage[UTF8]{ctex}
\begin{document}

\maketitle

\section*{MILP1 用于固定 OR 和 PACU 容量}

求解器采用了第4.2.1节中论文中提出的原始MILP1模型。操作室和PACU的容纳能力$R_1$和$R_2$是固定输入参数，与论文中的描述完全一致。 

让$\mathcal{S} = \{1,2\}$表示阶段$1$为“OR”关系，而阶段$2$为“PACU”。对于阶段$s \in \mathcal{S}$，让$\mathcal{R}_s = \{1,\ldots,R_s\}$表示房间集合，并将时间范围划分为$\{0,\ldots,K\}$个离散的时间段。

\begin{equation}
\min \; \Pi =
\sum_{s \in \mathcal{S}} \sum_{r \in \mathcal{R}_s} C_s^o h_{sr}
+
\sum_{s \in \mathcal{S}} C_s^e R_s
\tag{1}
\end{equation}

\begin{equation}
\sum_{r \in \mathcal{R}_s} \sum_{t=0}^{K-p_{is}} x_{isr}^t = 1,
\qquad
\forall i \in \mathcal{I},\ \forall s \in \mathcal{S}
\tag{2}
\end{equation}

\begin{equation}
\sum_{i \in \mathcal{I}} \sum_{\tau=\max\{0,t-p_{is}\}}^{t-1} x_{isr}^{\tau} \le 1,
\qquad
\forall s \in \mathcal{S},\ \forall r \in \mathcal{R}_s,\ \forall t = 1,\ldots,K
\tag{3}
\end{equation}

\begin{equation}
\sum_{\tau=t}^{\min\{K-1,t+p_{is}-1\}} y_{isr}^{\tau} \ge p_{is} x_{isr}^t,
\qquad
\forall i \in \mathcal{I},\ \forall s \in \mathcal{S},\ \forall r \in \mathcal{R}_s,\ \forall t = 0,\ldots,K-p_{is}
\tag{4}
\end{equation}

\begin{equation}
\sum_{i \in \mathcal{I}} \sum_{r \in \mathcal{R}_s} y_{isr}^t = z_s^t,
\qquad
\forall s \in \mathcal{S},\ \forall t = 0,\ldots,K-1
\tag{5}
\end{equation}

\begin{equation}
\sum_{t=0}^{K-1} z_s^t = \sum_{i \in \mathcal{I}} p_{is},
\qquad
\forall s \in \mathcal{S}
\tag{6}
\end{equation}

\begin{equation}
\sum_{r \in \mathcal{R}_s} \sum_{t=0}^{K-p_{is}} (t+p_{is}) x_{isr}^t = f_{is},
\qquad
\forall i \in \mathcal{I},\ \forall s \in \mathcal{S}
\tag{7}
\end{equation}

\begin{equation}
\sum_{r \in \mathcal{R}_2} \sum_{t=0}^{K-p_{i2}} t \, x_{i2r}^t \ge f_{i1},
\qquad
\forall i \in \mathcal{I}
\tag{8}
\end{equation}

\begin{equation}
f_{i2} - (f_{i1} + p_{i2}) \le 0,
\qquad
\forall i \in \mathcal{I}
\tag{9}
\end{equation}

\begin{equation}
(t+1-g_{sr}) \le K\omega_{sr}^t,
\qquad
\forall s \in \mathcal{S},\ \forall r \in \mathcal{R}_s,\ \forall t = 0,\ldots,K-1
\tag{10}
\end{equation}

\begin{equation}
\sum_{i \in \mathcal{I}} y_{isr}^t \le K(1-\omega_{sr}^t),
\qquad
\forall s \in \mathcal{S},\ \forall r \in \mathcal{R}_s,\ \forall t = 0,\ldots,K-1
\tag{11}
\end{equation}

\begin{equation}
g_{sr} - T \le h_{sr},
\qquad
\forall s \in \mathcal{S},\ \forall r \in \mathcal{R}_s
\tag{12}
\end{equation}

\begin{equation}
z_s^t \le R_s,
\qquad
\forall s \in \mathcal{S},\ \forall t = 0,\ldots,K-1
\tag{13}
\end{equation}

\begin{equation}
x_{isr}^t,\ y_{isr}^t,\ \omega_{sr}^t \in \{0,1\},
\qquad
\forall i \in \mathcal{I},\ \forall s \in \mathcal{S},\ \forall r \in \mathcal{R}_s,\ \forall t = 0,\ldots,K-1
\tag{14}
\end{equation}

\begin{equation}
z_s^t,\ f_{is},\ g_{sr},\ h_{sr} \ge 0,
\qquad
\forall i \in \mathcal{I},\ \forall s \in \mathcal{S},\ \forall r \in \mathcal{R}_s,\ \forall t = 0,\ldots,K-1
\tag{15}
\end{equation}

约束条件(2)至(7)用于分配每位患者，避免重复，包括第一个被占用的位置。方程(3)使用了原始论文中的窗口$[\max(0,t-p_{is}),t-1]$，位于$t=1,\ldots,K$处（PDF第15页）。完成变量$f_{is}$等于实际完成时间；这使得方程(8)至(9)能够实施源文献中描述的无需等待的传输方式（PDF第14页至第15页）。


\paragraph{起始时间范围的说明：2026-10-23。}
每个$p_{is}$都代表一个正整数个插槽。 定义$x_{isr}^t=0$
对于$K-p_{is}<t\le K-1$；只有$0\le t\le K-p_{is}$是起始决策。
所有正式包含后续开始的求和运算，都遵循这种零的约定。
原始 PDF 文件中的第 15 页，在所有开始时间都适用其服务覆盖不等式（4）：如果开始时间较晚，那么剩余的二进制占用槽位数会少于 $p_{is}$，因此该不等式本身就迫使开始时间为零。
这种明确的约定，在将(4)式限制在“可完成服务的时间范围内”之后，能够恢复原有的边界条件。同时，这也与当前求解器起始变量的范围相匹配。
它仅从之前提取的模型中移除多余的、过时的记录；将这些记录设置为零后，所有实际的计划及其目标都得以保留。

\end{document}
